\documentclass[11pt,a4paper]{ctexart}

\usepackage[a4paper,margin=1in]{geometry}
\usepackage{amsmath,amssymb,amsthm,mathtools,booktabs}
\usepackage{enumitem}
\usepackage{graphicx}
\usepackage{tikz}
\usetikzlibrary{arrows.meta,positioning}
\usepackage[hidelinks]{hyperref}

\theoremstyle{definition}
\newtheorem{definition}{定义}
\newtheorem{proposition}{命题}
\theoremstyle{remark}
\newtheorem{remark}{注}

\newcommand{\cJ}{\mathcal{J}}
\newcommand{\cM}{\mathcal{M}}
\newcommand{\cV}{\mathcal{V}}
\newcommand{\LU}{\ell}
\newcommand{\yes}{$\checkmark$}
\newcommand{\no}{$\times$}

\title{带 AGV 运输与有限缓冲阻塞的柔性作业车间联合调度：\\建模、死锁规避与学习型构造策略}
\author{（作者）}
\date{2026-11-04（草稿）}

\begin{document}
\maketitle

\paragraph{草稿说明}
{\sloppy
本文是大学生创新训练项目（大创）论文初稿，对应仓库代码 \texttt{src/fjsp\_lab/logistics/}。
文中所有数值都取自 \texttt{experiments/results/} 下已归档的 6 个实验目录：
\texttt{l1\_benchmark}、\texttt{l1\_decoupled}、\texttt{l1\_large}、\texttt{l1\_search}、
\texttt{l2\_formal\_train}、\texttt{l2\_formal\_test}。
每个目录的 \texttt{manifest.json} 记录了对应的代码提交与输入文件哈希。
投稿前需要完成：\par}
\begin{itemize}[nosep]
  \item 参考文献中标注\textbf{[待核对]}或\textbf{[待补全]}的条目，按原文核对或补全作者、题目和卷期；只读过摘要的文献需读全文确认文中的表述；
  \item 投稿前再做一次数据库检索（WoS/Scopus）和引文追溯，确认第~\ref{sec:related} 节的定位；
  \item 按目标期刊或会议的模板调整格式与篇幅。
\end{itemize}

\begin{abstract}
在柔性作业车间中，机器加工与自动导引车（AGV）运输相互耦合；当机器前的缓冲容量有限、
加工完的工件必须等 AGV 取走才释放机器（阻塞）时，非延迟的派工规则还会陷入死锁。
已有的柔性作业车间运输调度（FJSPT）研究普遍假设缓冲无限，而考虑有限缓冲的学习型方法又不含运输。
本文在 FJSPT 的时间框架下（AGV 有位置，空驶与载货行驶分别计时，工件从装卸站出发）引入有限输入缓冲、
加工后阻塞和派车预留，建立以最大完工时间为目标的约束模型，给出 CP-SAT 参考模型与独立的合法性检查器，
并在公开基准上复现了全部 10 个参考解。在此基础上提出前瞻式死锁规避：每个决策先在仿真副本上用保守策略推演到结束，
推演成功才执行。它既用来包装派工规则，也作为强化学习的动作掩码。学习型策略为每个“工序—机器—AGV”候选动作打分，
先模仿最好的规避规则，再用带共享基线的 REINFORCE 微调。实验表明：
（1）无缓冲加阻塞使小算例的最优值平均上升 6.8\%--12.8\%，使大算例的启发式解上升约 40\%--50\%；
（2）先排机器、再排车辆的解耦方法比联合优化平均差 5\%--10\%，无缓冲时升至 17\%--20\%，且有 2--3 个算例的机器方案无法执行；
（3）在预先登记的检验中，学习型策略在 3 个训练种子上都显著优于规则和同等时间的模拟退火；
在比训练算例大得多的 Brandimarte 算例上，每例采样 32 次取最好，比逐例最好规则好 16.3\%，比大预算模拟退火好 9.5\%；
但在小算例上距证得最优仍有约 12.5\% 的差距。

\noindent\textbf{关键词}：柔性作业车间调度；AGV；有限缓冲；阻塞；死锁规避；强化学习
\end{abstract}

\begin{center}
\begin{minipage}{0.9\textwidth}
\small\textbf{Abstract.}
We study flexible job shop scheduling with AGV transport (FJSPT) under finite input buffers and blocking after processing,
where non-delay dispatching may deadlock. We give a constraint model with explicit empty and loaded travel, a CP-SAT reference
model and an independent feasibility checker that replays all ten published benchmark solutions. A lookahead deadlock guard,
which accepts a decision only if a conservative rollout on a copy of the simulator finishes every job, wraps dispatching rules
and serves as the action mask of a learned per-candidate scoring policy trained by imitation and REINFORCE with a shared baseline.
Blocking without buffers raises small-instance optima by 6.8\%--12.8\% and large-instance heuristic makespans by about 40\%--50\%;
a decoupled machine-then-vehicle approach is 5\%--10\% worse than joint optimization, 17\%--20\% without buffers.
In pre-registered tests over three training seeds the policy beats the rules and equal-time simulated annealing; on the larger
Brandimarte instances 32 samples improve on the best rule by 16.3\% and on long-run annealing by 9.5\%, while on small instances
the gap to proven optima is still about 12.5\%.
\end{minipage}
\end{center}

\section{引言}

柔性作业车间调度问题（FJSP）要为每道工序选择机器并确定加工顺序，是生产调度中研究最多的问题之一\cite{dauzere2024}。
在实际车间里，工件要由 AGV 在装卸站与机器之间搬运。运输时间不可忽略，AGV 卸货后还要空驶到下一个取货点，
空驶时间取决于前后两次任务的位置，所以机器调度和车辆调度必须一起考虑。
这类问题称为带运输资源的柔性作业车间调度（FJSPT），从 Bilge 与 Ulusoy 的工作\cite{bilge1995}开始已有大量研究，
近年的精确与元启发式方法见文献\cite{homayouni2021,berterottiere2024}。

现有 FJSPT 研究，尤其是深度强化学习方法\cite{link2026,hadqn2025,wang2024}，普遍假设机器前后的缓冲无限：
工件随时可以送到机器前排队，加工完立即离开机器。实际车间的工位空间有限，常见的情形是机器前只有一两个等待位，
加工完的工件留在机器上，直到 AGV 取走才释放机器。这时会出现两个新问题：
\begin{enumerate}[nosep]
  \item \textbf{缓冲和阻塞改变了问题的结构。}机器被成品占着时不能开工，AGV 来晚了会直接延误下游；
        机器分配、加工顺序和运输顺序之间的耦合更强。
  \item \textbf{死锁。}机器 A 上的成品要去机器 B，B 上的成品要去 A，而两台机器前都没有空位，两个工件就会永远等下去。
        按非延迟规则派工时，这种情形并不罕见：本文的大算例上，缓冲容量不超过 2 时普通规则经常死锁，无缓冲时几乎全部死锁（第~\ref{sec:deadlock} 节）。
\end{enumerate}

有限缓冲与死锁在柔性制造系统的 Petri 网研究中有很长的历史，Hu 等\cite{hu2020}用 Petri 网加深度 Q 网络处理了有限资源与死锁避免；
Zhang 等\cite{zhang2026}用强化学习求解有限缓冲的柔性作业车间，但不含运输；也有中文期刊论文在流水车间类问题中把运输、阻塞与强化学习结合起来\cite{wangdh2026,yunn2026}。
据我们的检索（见第~\ref{sec:related} 节），同时具备“柔性路线、AGV 显式行驶（含空驶）、有限缓冲与阻塞、死锁处理、学习型方法”的工作尚未见报道。
柔性路线使机器选择和运输目的地都在决策时才确定，这与固定路线的流水车间不同。

本文的主要工作如下：
\begin{enumerate}[nosep]
  \item \textbf{模型与参考解。}在 FJSPT 的时间框架下引入有限输入缓冲、加工后阻塞和派车预留，给出约束模型、
        CP-SAT 参考模型与独立的合法性检查器。按基准约定，公开的 10 个参考解全部复现；CP-SAT 在 9 个算例上证得与参考值相同的最优值，
        在 fjsp7 上找到 108，优于公开的 112。
  \item \textbf{缓冲与协同的定量分析。}在小算例上用证得最优值、在大算例上用规则与模拟退火，量化阻塞与缓冲容量的影响；
        并在同样的预算下比较联合优化与“先排机器、再排车辆”的解耦优化，得到解耦代价随缓冲变紧的变化。
  \item \textbf{前瞻式死锁规避与学习型策略。}提出基于保守推演的死锁规避，把它同时用于规则基线和强化学习的动作掩码；
        学习型策略对“工序—机器—AGV”联合动作逐个打分，只有 704 个参数，先模仿再用 REINFORCE 微调。
  \item \textbf{预先登记的评价。}训练前固定假设、预算与比较口径，3 个训练种子分别报告；测试集只使用一次。
\end{enumerate}

第 2 节回顾相关工作；第 3 节给出问题与模型；第 4 节介绍求解方法；第 5 节报告实验；第 6 节总结并讨论局限。

\section{相关工作}\label{sec:related}

\paragraph{机器与运输联合调度。}
Bilge 与 Ulusoy\cite{bilge1995}研究了柔性制造系统中机器与物料搬运系统的同时调度，其布局和算例被后来的研究反复使用。
Deroussi 与 Norre\cite{deroussi2010}把它推广到柔性作业车间，给出了本文使用的 10 个小算例。
Homayouni 与 Fontes\cite{homayouni2021}提出了生产与运输联合调度的模型和算法，并给出了大算例所用的布局；
Berterottière 等\cite{berterottiere2024}给出了带运输资源的柔性作业车间的模型与方法，并公开了参考解。
Zhuang 等\cite{zhuang2025}用混合整数规划求解分区柔性制造系统中机器与 AGV 的同时调度，其中加工区不设缓冲；
Moinuddin 等\cite{moinuddin2026}用约束规划处理异构运输机器人与分区，并用路径预订避免碰撞。
这些工作或假设缓冲充足，或采用精确与元启发式方法，没有使用学习型方法。

\paragraph{有限缓冲、阻塞与死锁。}
在有限缓冲和阻塞下，调度必须避免循环等待。Petri 网是柔性制造系统死锁控制的主流工具，
Hu 等\cite{hu2020}在 Petri 网上用图卷积与深度 Q 网络做动态调度，并处理死锁避免。
中文文献中也有柔性作业车间在有限缓冲与 AGV 下的遗传算法研究，以及用时序 Petri 网求解机器人柔性作业车间无死锁调度的工作（见查新报告）。
本文不建 Petri 网，而是直接在事件仿真器上做前瞻推演，这样可以沿用 FJSPT 的时间模型和公开基准，与该领域的精确解直接比较。

\paragraph{柔性作业车间的深度强化学习。}
Zhang 等\cite{zhang2020l2d}在作业车间中用图神经网络和策略梯度学习派工，Song 等\cite{song2023}把这一思路推广到柔性作业车间；
这一线的方法通常把“工序—机器”对作为动作。对于带运输的问题，Link 等\cite{link2026}比较了联合学习与模块化学习的“协调差距”，
但其模型是固定路线的作业车间且缓冲无限；HA-DQN\cite{hadqn2025}和 Wang 等\cite{wang2024}研究了带运输约束的柔性作业车间，同样假设缓冲无限。
有限缓冲的学习型方法\cite{zhang2026}没有运输；运输、阻塞与强化学习的结合见于流水车间类问题\cite{wangdh2026,yunn2026}。
训练方法上，本文采用 REINFORCE\cite{williams1992}，基线取同一实例多次采样的平均值，与 Kool 等\cite{kool2019}、Kwon 等\cite{kwon2020}的共享基线思路相同。

\paragraph{本文的定位。}
表~\ref{tab:position}对比了最接近的工作。本文的贡献是组合式的：把学习型构造策略推广到带 AGV 显式行驶、
有限缓冲与阻塞的柔性作业车间，并用前瞻推演保证死锁规避；同时给出缓冲容量与联合/解耦差距之间关系的定量结果。

\begin{table}[ht]
\centering
\caption{与最接近工作的对比（\yes\ 有，\no\ 无，? 未能核实）}
\label{tab:position}
\small
\resizebox{\textwidth}{!}{%
\begin{tabular}{lccccc}
\toprule
文献 & 柔性路线 & AGV 显式行驶 & 有限缓冲/阻塞 & 死锁处理 & 方法 \\
\midrule
Hu 等 2020\cite{hu2020} & \yes & ? & \yes / ? & \yes & Petri 网 + DQN \\
Zhang 等 2026\cite{zhang2026} & \yes & \no & \yes / \no & \no & PPO + 异构图 \\
Link 等 2026\cite{link2026} & \no & \yes & \no / \no & \no & MAPPO \\
HA-DQN 2025\cite{hadqn2025} & \yes & \yes & \no / \no & \no & DQN \\
Zhuang 等 2025\cite{zhuang2025} & \yes & \yes & 无缓冲 / \yes & 假设不发生 & MILP \\
Moinuddin 等 2026\cite{moinuddin2026} & \yes & \yes & \no / \no & 路径预订 & CP \\
本文 & \yes & \yes & \yes / \yes & 前瞻推演 & 模仿 + REINFORCE \\
\bottomrule
\end{tabular}}
\end{table}

\section{问题描述与数学模型}

\subsection{问题描述与假设}

车间有机器集合 $\cM=\{1,\dots,m\}$ 和一个装卸站 $\LU$，工件集合 $\cJ$。工件 $j$ 有 $n_j$ 道按顺序加工的工序，
工序 $(j,o)$ 可在机器集合 $\cM_{jo}\subseteq\cM$ 中任选一台，在机器 $i$ 上的加工时间为 $p_{joi}$。
$V$ 台同质、单载 AGV 组成集合 $\cV$，0 时刻都停在 $\LU$。地点之间的行驶时间为 $\tau_{ab}$（$a,b\in\cM\cup\{\LU\}$），
空驶与载货行驶使用同一矩阵，对角线为 0；每次载货行驶另加装卸时间 $h\ge 0$。
机器 $i$ 前有容量为 $B_i\in\mathbb{Z}_{\ge 0}\cup\{\infty\}$ 的输入缓冲。其余假设：
\begin{itemize}[nosep]
  \item 工件 0 时刻都在 $\LU$；每道工序加工一次，不可中断；运输一旦开始不可中断。
  \item 同一工件在同一台机器上连续加工的工序之间不需要运输。
  \item \textbf{阻塞}：加工完的工件留在机器上，直到 AGV 取走才释放机器；机器被占着时不能开始其他工件的加工。
  \item \textbf{工位容量}：机器 $i$ 的工位共 $1+B_i$ 个位置（机器本身加缓冲）。派车时即在目的工位预留一个位置，
        工件被取走（或最后一道工序结束且不需运走）时释放。$\LU$ 容量无限。
  \item 两个开关：阻塞开关 $\beta\in\{0,1\}$，关闭时工件加工完立即进入容量无限的输出区，不占机器也不占工位；
        送回开关 $R\in\{0,1\}$，打开时工件最后要运回 $\LU$。FJSPT 基准算例的约定是 $\beta=0$、$R=0$、$B_i=\infty$、$h=0$。
  \item 不考虑路径冲突、拥堵和充电。
\end{itemize}
目标是最小化最大完工时间 $C_{\max}$。$R=1$ 时工件完工时刻为运回 $\LU$ 的送达时刻，否则为最后一道工序的结束时刻。

\subsection{约束模型}

工件 $j$ 的运输以“送往的工序号”$o$ 标识，$o=0,\dots,n_j-1$，$R=1$ 时另有送回 $\LU$ 的运输 $o=n_j$。
记 $\mu_{jo}$ 为工序 $(j,o)$ 所选机器，运输 $(j,o)$ 的起点 $a_{jo}=\mu_{j,o-1}$（$o=0$ 时为 $\LU$），
终点 $b_{jo}=\mu_{jo}$（$o=n_j$ 时为 $\LU$）。决策变量：
\begin{itemize}[nosep]
  \item $x_{joi}\in\{0,1\}$：工序 $(j,o)$ 是否在机器 $i$ 上加工；$S_{jo},C_{jo}$：开工与完工时刻；
  \item $y_{jo}\in\{0,1\}$：运输 $(j,o)$ 是否需要；$z_{jov}\in\{0,1\}$：是否由 AGV $v$ 执行；
  \item $D_{jo}\le P_{jo}\le Q_{jo}$：派车、取货、送达时刻。
\end{itemize}

\noindent\textbf{加工}：
\begin{align}
& \textstyle\sum_{i\in\cM_{jo}} x_{joi}=1,\qquad C_{jo}=S_{jo}+\sum_{i} p_{joi}x_{joi}, \label{eq:assign}\\
& S_{jo}\ge C_{j,o-1}\quad (o\ge 1). \label{eq:prec}
\end{align}

\noindent\textbf{运输}：$y_{j0}=1$，$R=1$ 时 $y_{jn_j}=1$；$1\le o<n_j$ 时 $y_{jo}=1$ 当且仅当 $\mu_{jo}\ne\mu_{j,o-1}$。
\begin{align}
& Q_{jo}=P_{jo}+y_{jo}\,(\tau_{a_{jo}b_{jo}}+h), \label{eq:trip}\\
& P_{jo}\ge C_{j,o-1}\ (o\ge 1),\qquad S_{jo}\ge Q_{jo}\ (o<n_j), \label{eq:handover}\\
& y_{jo}=0\ \Rightarrow\ D_{jo}=P_{jo}=C_{j,o-1}. \label{eq:notrip}
\end{align}

\noindent\textbf{车辆}：$\sum_{v} z_{jov}=y_{jo}$。同一 AGV 的运输串成一条路径：若运输 $k'$ 紧接在 $k$ 之后，则
\begin{equation}
D_{k'}\ge Q_k,\qquad P_{k'}\ge D_{k'}+\tau_{b_k a_{k'}}; \label{eq:vehicle}
\end{equation}
若 $k$ 是该车的第一次运输，则 $D_k\ge 0$、$P_k\ge D_k+\tau_{\LU a_k}$。

\noindent\textbf{阻塞}：工序 $(j,o)$ 对机器的占用区间为 $[S_{jo},F_{jo})$，其中
\begin{equation}
F_{jo}=\begin{cases}
P_{j,o+1}, & \beta=1 \text{ 且 } y_{j,o+1}=1,\\
S_{j,o+1}, & y_{j,o+1}=0 \text{（工件留在本机加工下一道工序）},\\
C_{jo}, & \beta=0,\ \text{或 } o=n_j-1 \text{ 且 } R=0.
\end{cases}\label{eq:release}
\end{equation}
同一机器上不同工序的占用区间互不重叠：若 $x_{joi}=x_{j'o'i}=1$，则 $F_{jo}\le S_{j'o'}$ 或 $F_{j'o'}\le S_{jo}$。
$\beta=1$ 时这比“加工不重叠”更严格。

\noindent\textbf{工位容量}：工序 $(j,o)$ 在其机器的工位上从 $r_{jo}$ 起占用一个位置，其中 $y_{jo}=1$ 时 $r_{jo}=D_{jo}$（派车即预留），
否则 $r_{jo}=S_{jo}$；释放时刻 $\beta=1$ 时同式~\eqref{eq:release}，$\beta=0$ 时为 $C_{jo}$（留在本机时为 $S_{j,o+1}$）。对任意机器 $i$ 和时刻 $t$，
\begin{equation}
\bigl|\{(j,o):x_{joi}=1,\ r_{jo}\le t<\text{释放时刻}\}\bigr|\le 1+B_i. \label{eq:buffer}
\end{equation}

\noindent\textbf{目标}：
\begin{equation}
\min C_{\max},\qquad C_{\max}\ge \begin{cases} Q_{jn_j}, & R=1,\\ C_{j,n_j-1}, & R=0.\end{cases}
\end{equation}

所有时间取整数，区间左闭右开。式~\eqref{eq:vehicle} 使空驶时间依赖前后两次运输的位置，而位置又取决于机器选择，
这是机器与车辆必须联合决策的根本原因。

\begin{proposition}\label{prop:feasible}
对任意算例和任意 $B_i\ge 0$，可行解都存在。
\end{proposition}
\begin{proof}
构造串行调度：工件逐个处理，前一个工件完工（$R=1$ 时运回 $\LU$）后才派出下一个；只用一台 AGV，
每次运输都等上一道工序结束再派车，每道工序任选一台可选机器。任一时刻车间里最多一个工件，
所以各工位最多被占一个位置（$\le 1+B_i$），机器占用区间不重叠，车辆路径的时刻也都满足式~\eqref{eq:vehicle}。
\end{proof}

\begin{remark}
取 $\tau\equiv 0$、$h=0$、$\beta=0$、$R=0$、$B_i=\infty$ 时，运输不占时间，车辆约束不再起作用，模型退化为 FJSP；作业车间问题已是 NP 难的\cite{garey1976}，
所以本问题也是 NP 难的。命题~\ref{prop:feasible} 说明死锁不是问题本身无解，而是派工方式造成的。
\end{remark}

\subsection{死锁}

图~\ref{fig:deadlock} 是 $B=0$ 时的死锁：工件 1 在 $M_1$ 上加工完，下一道工序只能在 $M_2$ 上做；
工件 2 在 $M_2$ 上加工完，下一道工序只能在 $M_1$ 上做。两台机器都被成品占着，没有空位可以预留，两个工件都无法派车。
模型允许两台 AGV 在同一时刻分别取走两个工件并互换机器（取货时机器即空出），但逐个决策的仿真器做不到这种同时互换。
非延迟规则只看当前，一旦走进这样的状态就无法退出。

\begin{figure}[ht]
\centering
\begin{tikzpicture}[box/.style={draw,minimum width=2.4cm,minimum height=1cm,align=center},>=Stealth]
  \node[box] (m1) {$M_1$（$B=0$）\\工件 1 已加工完};
  \node[box,right=3.2cm of m1] (m2) {$M_2$（$B=0$）\\工件 2 已加工完};
  \draw[->,thick] (m1.north east) to[bend left=20] node[above,font=\small]{工件 1 要去 $M_2$，无空位} (m2.north west);
  \draw[->,thick] (m2.south west) to[bend left=20] node[below,font=\small]{工件 2 要去 $M_1$，无空位} (m1.south east);
\end{tikzpicture}
\caption{阻塞且无缓冲时的循环等待}
\label{fig:deadlock}
\end{figure}

\subsection{CP-SAT 参考模型与合法性检查器}

小规模算例用 OR-Tools CP-SAT\cite{ortools}求解上述模型：每道工序在每台可选机器上一个可选区间，恰好一个出现；
每台 AGV 的运输用一个 \texttt{AddCircuit} 约束串成回路（车场即 0 时刻的 $\LU$）；阻塞用 \texttt{NoOverlap}，
工位容量用 \texttt{Cumulative}；行驶时间依赖机器选择，用 \texttt{Element} 约束取值。
求解器状态为 OPTIMAL 时才称最优，FEASIBLE 只作为限时可行解。

所有方法返回的调度都必须通过独立的合法性检查器。检查器把生产部分交给不改动的静态 FJSP 检查器
（每道工序恰好一次、工艺顺序、可选机器、加工时长、加工不重叠），再逐条检查：
所需运输恰好各一次且起讫正确；运输时长；取货不早于工序结束、开工不早于送达；
车辆从 $\LU$ 出发、运输串联且空驶时间足够；阻塞下占用区间不重叠；工位占用不超过 $1+B_i$；
报告的 $C_{\max}$ 等于独立重算值。检查器只判定，不修复。

\section{求解方法}

\subsection{事件驱动仿真与候选动作}

规则和学习型策略都在一个事件驱动仿真器上构造调度。仿真器在“有事可做”的时刻（非延迟）停下来，给出候选动作集合。
一个候选动作是三元组决策：把等待中的工件 $j$ 用空闲 AGV $v$ 送往下一道工序的某台可选机器 $i$（或送回 $\LU$）；
若工件所在机器也能加工下一道工序，还可以选择“原地留下”。目的工位已满的候选不出现。
派车时锁定机器并预留目的工位；工件到达后在输入缓冲中按到达顺序等待，机器空闲时开工。
同一时刻的事件按固定顺序处理（工序结束、取货、送达，然后开工，最后决策），保证结果可复现。

\subsection{前瞻式死锁规避}

对一个候选动作 $c$，在仿真器副本上执行 $c$，再用\textbf{保守推演策略}运行到结束：
已进入车间的工件按联合最早完工（EFT）推进，车间清空之后才从 $\LU$ 放行下一个工件。
推演完成则称 $c$ \textbf{安全}。推演成功即证明执行 $c$ 后的状态仍能完成全部工件；推演失败不一定表示真的会死锁，
所以这种规避是保守的。

规避包装器按策略的偏好顺序逐个检查候选，执行第一个安全的；都不安全则等待下一个事件。
缓冲全部无限时不可能死锁（等待的工件总能被运走），不做检查，此时包装后的规则与原规则完全相同。
在实验中，带规避的规则 420 次运行、学习型策略的全部训练与测试回合都没有出现死锁（第~\ref{sec:deadlock}、\ref{sec:l2} 节）。
从推演的构造看，初始状态总是安全的，规避后每一步也保持安全，这是规避有效的直观原因；
但同一时刻事件的处理顺序会让“执行后继续推演”与“原推演”略有不同，我们没有给出严格的完备性证明。

\subsection{规则基线}

联合规则 EFT 在全部“工序—机器—AGV”候选中选预计完工时刻最早者，同时考虑空驶、载货行驶和目的机器的排队。
解耦规则先按 FIFO、SPT 或 MWKR 选工件，并在不计运输的前提下为它选机器，再派最近的空闲车辆（记为 +NV），
对应“先排生产、再派车辆”的做法。下文的“规避规则”指加了前瞻规避的 EFT、FIFO+NV、SPT+NV、MWKR+NV，
“逐例最好规则”指这 4 条中在该实例上 $C_{\max}$ 最小者。

\subsection{学习型构造策略}

\paragraph{决策过程。}状态为仿真器在决策时刻的全部信息，动作为一个候选，只在结束时给一次回报 $-C_{\max}/S$，
其中尺度 $S=\sum_{j,o}\min_i p_{joi}/m$。

\paragraph{候选特征与打分。}每个候选用 20 维特征描述（表~\ref{tab:features}），时间量都除以 $S$。
EFT、SPT、MWKR 和最近车辆等规则用到的信息都能由这些特征表示。打分函数为两层感知机
\begin{equation}
f(\boldsymbol{x})=\boldsymbol{w}_2^{\top}\tanh(W_1\boldsymbol{x}+\boldsymbol{b}_1),
\end{equation}
隐藏层 32 个单元，共 704 个参数。设 $A$ 为当前候选集合、$U\subseteq A$ 为尚未被屏蔽的候选，策略为
\begin{equation}
\pi(c\mid s,U)=\frac{\exp f(\boldsymbol{x}_c)}{\sum_{c'\in U}\exp f(\boldsymbol{x}_{c'})},\qquad c\in U.
\end{equation}

\begin{table}[ht]
\centering
\caption{候选特征（20 维）}
\label{tab:features}
\small
\begin{tabular}{ll}
\toprule
类别 & 特征 \\
\midrule
候选本身 & 是否原地留下、是否送回 $\LU$、空驶时间、载货行驶时间、到达后等机器的时间、\\
         & 加工时间、预计完工时刻、预计完工时刻在候选中的排名 \\
工件 & 剩余最小工作量、剩余工序数、已等待时间 \\
目的工位 & 占用率、本步后是否占满、机器是否被占着、输入缓冲排队数 \\
出发工位与系统 & 出发机器的排队数、空闲车辆比例、已完工比例、是否阻塞、缓冲是否有限 \\
\bottomrule
\end{tabular}
\end{table}

\paragraph{死锁规避掩码。}从 $U=A$ 开始，按 $\pi$ 选出候选（贪心时取概率最大者，采样时随机抽取），用前瞻推演检查；
不安全就把它移出 $U$ 并重新归一化再选；$U$ 为空则等待。这样采样得到的动作恰好服从 $\pi$ 限制在安全候选上的分布，
而且只检查被选中的候选，比先检查全部候选省时。策略梯度在接受该动作时的 $U$ 上计算，被屏蔽的候选不接收梯度。
缓冲无限时不做检查。

\paragraph{训练。}分两个阶段：
\begin{enumerate}[nosep]
  \item \textbf{模仿}：对每个训练实例，记录逐例最好规则的每一步决策，用交叉熵训练打分函数。
  \item \textbf{REINFORCE}：对每个实例采样 $K$ 条轨迹，基线取这 $K$ 次的平均完工时间 $b$，第 $k$ 条轨迹的优势为
        $A_k=(b-C_{\max}^{(k)})/b$，按
        \begin{equation}
        \nabla_\theta J\approx\frac{1}{K}\sum_{k=1}^{K}A_k\sum_{t}\nabla_\theta\log\pi_\theta\bigl(c_t^{(k)}\mid s_t^{(k)},U_t^{(k)}\bigr)
        \end{equation}
        用 Adam 更新\cite{williams1992,kool2019,kwon2020}。
\end{enumerate}
只用验证集选模型，指标为“贪心 $C_{\max}$ / 逐例最好规则”的平均值。

\paragraph{推理。}贪心推理构造一次调度；采样推理在贪心解之外再采样 $N$ 次（本文 $N=8,32$），取最好者。
学习到的策略本身不保证可行性，可行性由仿真器、死锁规避与检查器共同保证。

\subsection{对照方法}

\paragraph{模拟退火（SA）。}编码分四段：工序顺序、机器分配、运输顺序、车辆分配。
解码先按给定顺序用拓扑排序求最早时间表，缓冲容量作为前后约束；顺序前后矛盾时，先把运输改为按工序顺序排列再试，
仍矛盾则退回优先级解码（仿真加死锁规避）。邻域为工序或运输顺序的交换与插入，以及机器或车辆的重新分配。
SA 从 4 条规避规则的最好解出发，温度按起始解的比例几何下降，所以结果不劣于规避规则。
在同样的编码上我们也实现了遗传算法（GA），在小算例上总体不优于 SA，下文以 SA 为主要参考。

\paragraph{解耦优化。}第一阶段只排机器：D0 不计运输；D1 把载货行驶时间当作工序间隔，假设车辆随时可用。
第二阶段固定机器方案（各机器的工序及顺序），用 CP-SAT 只优化车辆与时间。机器方案无法执行时如实记录，不做修补。

\section{实验}\label{sec:exp}

\subsection{算例、配置与实现}

\begin{itemize}[nosep]
  \item \textbf{小算例}：Deroussi 与 Norre\cite{deroussi2010}的 10 个 FJSPT 算例（fjsp1--10，5--8 个工件，8 台机器）及其布局，
        参考解来自 Berterottière 等\cite{berterottiere2024}公开的数据。
  \item \textbf{大算例}：Brandimarte\cite{brandimarte1993}的 mk01--05、mk07（10--20 个工件，55--150 道工序，4--8 台机器），
        行驶时间取 Homayouni 与 Fontes\cite{homayouni2021}的 4/5/6/8 机整数布局。mk06、mk08--10 所需的 10 机、15 机布局含小数时间，未纳入。
  \item \textbf{配置}：基准约定（不阻塞、缓冲无限），以及“阻塞 + 输入缓冲 $B=\infty/2/1/0$”，共 5 种；都是 2 台 AGV、不送回 $\LU$。
        大算例的车辆数是我们设定的，数值不能与原文献直接比较。
  \item \textbf{实现}：Python 3.11、NumPy 2.4、OR-Tools 9.15。CP-SAT 每例 60 s、8 线程；学习型策略只用 CPU。
\end{itemize}

\subsection{基准对齐与参考解}

按基准约定，10 个公开参考解重放后 $C_{\max}$ 全部等于作者报告值，并都通过检查器；
若改为阻塞，或把装卸站放在最后一个地点，则不一致。这说明本文的模型语义与 FJSPT 基准一致。
CP-SAT 在 fjsp1--6、8--10 上证得与参考值相同的最优值（134、114、120、114、94、138、178、144、174）。
在 fjsp7 上，CP-SAT 于限时内找到 108（下界 104）；同一算例在“阻塞 + 无限缓冲”配置下证得最优值为 108，
而阻塞配置的可行解在基准约定下同样可行，所以基准约定下的最优值在 $[104,108]$ 内，优于公开参考值 112。

\subsection{阻塞与缓冲容量的影响}

表~\ref{tab:buffer} 给出各配置相对基准约定的 $C_{\max}$ 平均增幅。小算例上用 CP-SAT，未证得最优时用解与下界给出区间；
大算例上 CP-SAT 在 60 s 内找不到解（从 55 道工序起，阻塞并限制缓冲后即如此），改用最好的规避规则和 SA。

\begin{table}[ht]
\centering
\caption{相对基准约定的 $C_{\max}$ 平均增幅}
\label{tab:buffer}
\small
\resizebox{\textwidth}{!}{%
\begin{tabular}{lcccc}
\toprule
 & 阻塞 + $B=\infty$ & 阻塞 + $B=2$ & 阻塞 + $B=1$ & 阻塞 + $B=0$ \\
\midrule
小算例，CP-SAT（证得最优数） & 0.21\%--0.59\%（10） & 0.21\%--0.59\%（10） & $-$0.16\%--1.17\%（9） & 6.8\%--12.8\%（5） \\
大算例，最好规避规则 & 4.1\% & 6.0\% & 11.5\% & 51.8\% \\
大算例，SA（2 个种子取优） & 4.2\% & 3.0\% & 9.0\% & 40.7\% \\
\bottomrule
\end{tabular}}
\end{table}

三种方法的定性结论一致：阻塞本身使 $C_{\max}$ 上升约 4\%（大算例），缓冲降到 1 时约 10\%，无缓冲时大算例上升约 40\%--50\%。
小算例只有 5--8 个工件，缓冲位很少被占满，$B\ge1$ 时影响不到约 1\%，所以缓冲的作用要在更大的算例上研究。
更强的方法（SA）把 $B=0$ 的代价估计从 52\% 降到 41\%，说明其中一部分来自规则本身太弱。

\subsection{联合优化与解耦优化的差距}

表~\ref{tab:decoupled} 在小算例上比较联合 CP-SAT 与两种解耦方法，总预算相同（解耦的两个阶段各 30 s）。
可执行的情况下解耦方法的两个阶段都证得最优，所以差距不是限时造成的。

\begin{table}[ht]
\centering
\caption{协调差距 =（解耦 $-$ 联合）/ 联合，小算例}
\label{tab:decoupled}
\small
\begin{tabular}{lccc}
\toprule
配置 & D0 平均（最大） & D1 平均（最大） & 机器方案无法执行 \\
\midrule
基准约定 & 10.1\%（17.5\%） & 5.5\%（11.9\%） & 0 / 0 \\
阻塞 + $B=\infty$ & 11.2\%（22.8\%） & 5.1\%（9.3\%） & 0 / 0 \\
阻塞 + $B=2$ & 9.7\%（18.5\%） & 4.8\%（11.1\%） & 0 / 0 \\
阻塞 + $B=1$ & 9.8\%（22.8\%） & 4.5\%（7.4\%） & 0 / 0 \\
阻塞 + $B=0$ & 20.1\%（34.7\%） & 17.3\%（28.8\%） & D0 2 / D1 3 \\
\bottomrule
\end{tabular}
\end{table}

在 95 个可执行的情况里，联合 CP-SAT 每一个都严格更好。只要有至少 1 个缓冲位，解耦代价大致保持在 D1 约 5\%、D0 约 10\%；
阻塞且无缓冲时代价升到 17\%--20\%，还有 2--3 个算例的机器方案根本无法执行，而联合 CP-SAT 在这些算例上都找到了可行调度。
即使第一阶段计入了运输时间（D1），由于不考虑阻塞，也照样会排出无法执行的方案。
联合 CP-SAT 在少数配置下只有限时可行解，此时真实差距只会更大，表中数字偏保守。

值得注意的是，简单的联合贪心规则 EFT 在小算例上并不比解耦规则好（无缓冲时与逐例最好解耦规则相比 0 胜 10 负），
说明联合决策的价值要靠更强的优化能力才能发挥，这也是本文引入学习型方法的动机。

\subsection{大算例上的死锁与规避}\label{sec:deadlock}

表~\ref{tab:deadlock} 给出普通（不加规避）规则在大算例上的死锁次数。$B\le2$ 时死锁已经很常见，$B=0$ 时几乎全部死锁，
所以在这类问题上，不处理死锁的非延迟规则不能作为基线。加上前瞻规避后，420 次运行（6 算例 × 5 配置 × 14 个规则或随机种子）
全部给出合法调度，单次最长 1.2 s；不会死锁的情况下，结果与原规则完全相同。

\begin{table}[ht]
\centering
\caption{普通规则在大算例上的死锁次数 / 运行次数（RANDOM 为 10 个种子）}
\label{tab:deadlock}
\small
\begin{tabular}{lccccc}
\toprule
配置 & EFT & FIFO+NV & SPT+NV & MWKR+NV & RANDOM \\
\midrule
基准约定、阻塞 + $B=\infty$ & 0/6 & 0/6 & 0/6 & 0/6 & 0/60 \\
阻塞 + $B=2$ & 1/6 & 2/6 & 2/6 & 2/6 & 26/60 \\
阻塞 + $B=1$ & 3/6 & 5/6 & 3/6 & 6/6 & 45/60 \\
阻塞 + $B=0$ & 6/6 & 6/6 & 5/6 & 6/6 & 60/60 \\
\bottomrule
\end{tabular}
\end{table}

\subsection{学习型策略的评价}\label{sec:l2}

\paragraph{训练。}训练集与验证集由生成器产生，种子区间互不重叠：5--10 个工件，每个工件 2--5 道工序，每道工序 1--2 台可选机器，
加工时间 5--30；机器数 4/5/6/8，使用上述布局；2 台 AGV，5 种配置随机抽取。训练集 400 个、验证集 100 个。
模仿 15 轮（学习率 $3\times10^{-3}$），REINFORCE 120 轮（$K=8$，每批 8 个实例，学习率 $10^{-3}$）。
3 个训练种子（0、1、2）依次训练，每个种子用 18 个进程并行采样，各用时 29--34 分钟。
选中的检查点分别来自第 58、119、74 轮 REINFORCE，验证集上“贪心 / 逐例最好规则”为 1.001、1.002、0.998；
只做模仿时验证集上最好为 1.023--1.028。验证集同时用于选模型，这两组数字偏乐观，只作描述。

\paragraph{测试协议。}训练前固定以下设计并登记假设。测试集 T1 为小算例 10 个 × 5 配置（50 个组合），
T2 为大算例 6 个 × 5 配置（30 个组合）；测试集只在模型选定后使用一次。每个组合上运行：4 条规避规则；
每个种子的贪心、采样 8 次、采样 32 次；同等评估次数的 SA（33 次）；同等时间的 SA（时间上限取同一组合上采样 32 次的实测耗时）。
推理与同等时间 SA 在同一进程中顺序运行。另引用已归档的大预算 SA（2 个种子取优，T1 每次 5000 次评估、T2 每次 1000 次评估）
和 CP-SAT 结果作参照。假设采用单侧配对 Wilcoxon 符号秩检验，在每个种子内对 H1--H4 做 Holm 校正，3 个种子中至少 2 个满足即成立。

\paragraph{假设检验。}表~\ref{tab:hyp} 给出结果。4 个假设在 3 个种子上全部成立。H4 虽然显著，幅度只有约 2\%。

\begin{table}[ht]
\centering
\caption{预先登记的假设（胜/平/负按组合计数）}
\label{tab:hyp}
\small
\begin{tabular}{llcccc}
\toprule
假设 & 内容 & 种子 0 & 种子 1 & 种子 2 & Holm 校正 $p$ \\
\midrule
H1 & 贪心优于 EFT（T1$\cup$T2，$n=80$） & 58/2/20 & 56/4/20 & 61/2/17 & $\le 3.3\times10^{-6}$ \\
H2 & 采样 32 优于逐例最好规则（$n=80$） & 77/3/0 & 76/0/4 & 79/0/1 & $\le 4.9\times10^{-14}$ \\
H3 & 同等时间优于 SA（T2，$n=30$） & 28/0/2 & 28/0/2 & 27/0/3 & $\le 5.8\times10^{-6}$ \\
H4 & 同等时间优于 SA（T1，$n=50$） & 32/4/14 & 31/4/15 & 31/5/14 & $\le 5.7\times10^{-3}$ \\
\bottomrule
\end{tabular}
\end{table}

\paragraph{描述性结果。}表~\ref{tab:l2} 给出按配置的 $C_{\max}$ 比值（3 个种子平均，小于 1 表示策略更好）。

\begin{table}[ht]
\centering
\caption{学习型策略的 $C_{\max}$ 比值（3 个种子平均）}
\label{tab:l2}
\small
\resizebox{\textwidth}{!}{%
\begin{tabular}{llccccc}
\toprule
测试集 & 配置 & 贪心 / 最好规则 & 采样 32 / 最好规则 & 同时 SA / 最好规则 & 采样 32 / SA 参考 & 采样 32 / 最优（$n$） \\
\midrule
T1 & 基准约定 & 0.978 & 0.928 & 0.960 & 1.041 & 1.115（9） \\
   & $B=\infty$ & 1.000 & 0.937 & 0.959 & 1.032 & 1.135（10） \\
   & $B=2$ & 1.008 & 0.925 & 0.956 & 1.015 & 1.120（10） \\
   & $B=1$ & 0.999 & 0.921 & 0.958 & 1.028 & 1.122（9） \\
   & $B=0$ & 1.042 & 0.947 & 0.945 & 1.054 & 1.140（5） \\
   & 全部 & 1.005 & 0.932 & 0.956 & 1.034 & 1.125（43） \\
\midrule
T2 & 基准约定 & 0.843 & 0.797 & 0.981 & 0.833 & -- \\
   & $B=\infty$ & 0.891 & 0.818 & 0.981 & 0.865 & -- \\
   & $B=2$ & 0.914 & 0.814 & 0.980 & 0.880 & -- \\
   & $B=1$ & 0.925 & 0.844 & 0.971 & 0.902 & -- \\
   & $B=0$ & 1.079 & 0.911 & 0.935 & 1.044 & -- \\
   & 全部 & 0.931 & 0.837 & 0.970 & 0.905 & -- \\
\bottomrule
\end{tabular}}
\end{table}

\begin{itemize}[nosep]
  \item \textbf{大算例上优势最明显。}T2 的工件数和工序数都比训练算例大得多。贪心推理一次构造就比逐例最好规则好 6.9\%，
        与大预算 SA 持平（比值 1.010）；采样 32 次比逐例最好规则好 16.3\%，比大预算 SA 好 9.5\%，3 个种子分别在 26、24、25/30 个组合上更好。
        这说明策略学到的打分方式能推广到更大的问题上。
  \item \textbf{小算例上搜索方法仍然更强。}T1 上采样 32 次比大预算 SA 差 3.4\%，只在 5--7/50 个组合上更好；
        距 CP-SAT 证得最优平均 12.5\%（贪心 20.9\%），43 个证得最优的组合中没有一个达到最优值。
        作为比较，大预算 SA 距最优约 8\%--12\%。
  \item \textbf{$B=0$ 最弱（探索性）。}无缓冲时贪心推理在两组测试集上都不如逐例最好规则，T2 上采样 32 次后仍比大预算 SA 差 4.4\%；
        其余 4 种配置下，T2 的采样 32 次都优于大预算 SA（比值 0.83--0.90）。这与试训时的观察一致。
  \item \textbf{种子间差异很小。}3 个种子的胜负数和比值都很接近，结论不依赖某一个种子。
\end{itemize}

\paragraph{计算时间。}表~\ref{tab:time} 为单进程顺序运行的实测时间。大算例上有限缓冲配置的采样较慢：缓冲无限时不做前瞻推演，采样 32 次不到 1 s，所以多出的时间来自推演。

\begin{table}[ht]
\centering
\caption{推理时间（秒，单进程）}
\label{tab:time}
\small
\begin{tabular}{lccc}
\toprule
测试集 & 贪心平均 & 采样 32 平均 & 采样 32 最长 \\
\midrule
T1 & 0.006 & 0.20 & 0.43 \\
T2（全部配置） & 0.34 & 10.9 & 58.6 \\
T2（$B=2/1/0$） & -- & 16.1 / 17.1 / 19.8 & -- \\
\bottomrule
\end{tabular}
\end{table}

\subsection{讨论}

\begin{enumerate}[nosep]
  \item \textbf{缓冲是关键参数。}三种方法一致表明无缓冲加阻塞使完工时间大幅上升，而一两个缓冲位就能消除大部分损失；
        对车间设计而言，这提示在机器前保留少量等待位的价值。
  \item \textbf{联合决策的价值要靠优化能力兑现。}联合 CP-SAT 严格优于解耦优化，但联合贪心规则并不优于解耦规则。
        学习型策略在同一联合动作空间中学到了比规则更好的打分，尤其在大算例上。
  \item \textbf{学习型方法的适用范围。}它在大规模、短时间的情形下最有优势；在小算例或允许长时间搜索时，SA 和精确方法更好。
        这与毕设中维护调度问题上的观察一致。
  \item \textbf{无缓冲是难点。}$B=0$ 时前瞻规避最频繁地介入，可选动作少，策略与规则的差别也最小；针对这一配置的特征或训练比重可能是改进方向。
\end{enumerate}

\section{结论与展望}

本文研究带 AGV 运输、有限输入缓冲与加工后阻塞的柔性作业车间调度。在与公开基准完全对齐的模型上，
我们定量给出了阻塞与缓冲容量的影响，以及联合与解耦决策的差距随缓冲变紧的变化；
提出了以保守推演为基础的死锁规避，并把它作为学习型构造策略的动作掩码。
在预先登记的检验中，学习型策略显著优于规则与同等时间的模拟退火，在比训练算例更大的算例上优于大预算模拟退火。

\paragraph{局限。}
\begin{itemize}[nosep]
  \item 测试算例少（小算例 10 个、大算例 6 个），同一算例的 5 种配置之间不独立。
  \item 大算例没有精确参考，“大预算 SA”每次也只有 1000 次评估。
  \item “同等时间”的比较依赖两种 Python 实现的效率，换实现结论可能变化。
  \item 只考虑 2 台同质单载 AGV、静态问题，不含路径冲突与充电；死锁规避的完备性没有严格证明。
  \item 没有与文献中的深度强化学习方法直接比较，也没有做消融实验。
\end{itemize}

\paragraph{展望。}用学习型策略的解作为搜索的初始解，以弥补小算例上的差距；加快前瞻推演；
扩展到新工件到达、机器故障等动态事件，以及路径冲突与异构车辆。

\begin{thebibliography}{99}
\small

\bibitem{dauzere2024}
Dauzère-Pérès S, Ding J, Shen L, Tamssaouet K. The flexible job shop scheduling problem: a review.
\textit{European Journal of Operational Research}, 2024, 314(2): 409--432.

\bibitem{bilge1995}
Bilge Ü, Ulusoy G. A time window approach to simultaneous scheduling of machines and material handling system in an FMS.
\textit{Operations Research}, 1995, 43(6): 1058--1070.

\bibitem{deroussi2010}
Deroussi L, Norre S. Simultaneous scheduling of machines and vehicles for the flexible job shop problem.
International Conference on Metaheuristics and Nature Inspired Computing, 2010. \textbf{[待核对：会议名称与页码]}

\bibitem{homayouni2021}
Homayouni S M, Fontes D B M M. Production and transport scheduling in flexible job shop manufacturing systems.
\textit{Journal of Global Optimization}, 2021, 79(2): 463--502.

\bibitem{berterottiere2024}
Berterottière L, Dauzère-Pérès S, Yugma C. Flexible job-shop scheduling with transportation resources.
\textit{European Journal of Operational Research}, 2024, 312(3): 890--909.

\bibitem{brandimarte1993}
Brandimarte P. Routing and scheduling in a flexible job shop by tabu search.
\textit{Annals of Operations Research}, 1993, 41: 157--183.

\bibitem{zhuang2025}
Zhuang C, et al. Comprehensive MILP formulation and solution for simultaneous scheduling of machines and AGVs in a partitioned flexible manufacturing system.
\textit{Machines}, 2025, 13: 519. \textbf{[待核对：作者全名]}

\bibitem{moinuddin2026}
Moinuddin A, Er Raqabi E M, Van Hentenryck P. Scheduling and routing in the flexible job shop with heterogeneous transbots and zoning: a constraint programming approach.
arXiv:2604.24483, 2026.

\bibitem{link2026}
Link M, Hoss J, Klarmann N. An analysis of the coordination gap between joint and modular learning for job shop scheduling with transportation resources.
IEEE CASE 2026; arXiv:2604.24117.

\bibitem{hadqn2025}
A heuristic-assisted deep reinforcement learning algorithm for flexible job shop scheduling with transport constraints.
\textit{Complex \& Intelligent Systems}, 2025. doi:10.1007/s40747-025-01828-6. \textbf{[待核对：作者]}

\bibitem{wang2024}
Wang J, et al. Dynamic integrated scheduling of production equipment and automated guided vehicles in a flexible job shop based on deep reinforcement learning.
\textit{Processes}, 2024, 12(11): 2423. \textbf{[待核对：作者全名]}

\bibitem{hu2020}
Hu L, et al. Petri-net-based dynamic scheduling of flexible manufacturing system via deep reinforcement learning with graph convolutional network.
\textit{Journal of Manufacturing Systems}, 2020, 55. doi:10.1016/j.jmsy.2020.02.004. \textbf{[待核对：题目、作者与页码]}

\bibitem{zhang2026}
Zhang S, et al. arXiv:2602.24180, 2026. \textbf{[待补全：题目与作者；内容为带有限缓冲与配料约束的柔性作业车间强化学习调度]}

\bibitem{wangdh2026}
王东海, 等. 制造技术与机床, 2026. \textbf{[待补全：按知网补全题目与卷期；内容为带中央缓存的重入阻塞流水车间深度强化学习调度]}

\bibitem{yunn2026}
余娜娜, 等. 中国机械工程, 2026. \textbf{[待补全：按知网补全题目与卷期；内容为阻塞与异构 AGV 下的混合流水车间调度]}

\bibitem{zhang2020l2d}
Zhang C, Song W, Cao Z, Zhang J, Tan P S, Xu C. Learning to dispatch for job shop scheduling via deep reinforcement learning.
\textit{Advances in Neural Information Processing Systems}, 2020.

\bibitem{song2023}
Song W, Chen X, Li Q, Cao Z. Flexible job-shop scheduling via graph neural network and deep reinforcement learning.
\textit{IEEE Transactions on Industrial Informatics}, 2023, 19(2): 1600--1610. \textbf{[待核对]}

\bibitem{williams1992}
Williams R J. Simple statistical gradient-following algorithms for connectionist reinforcement learning.
\textit{Machine Learning}, 1992, 8: 229--256.

\bibitem{kool2019}
Kool W, van Hoof H, Welling M. Attention, learn to solve routing problems!
International Conference on Learning Representations, 2019.

\bibitem{kwon2020}
Kwon Y-D, Choo J, Kim B, Yoon I, Gwon Y, Min S. POMO: policy optimization with multiple optima for reinforcement learning.
\textit{Advances in Neural Information Processing Systems}, 2020.

\bibitem{garey1976}
Garey M R, Johnson D S, Sethi R. The complexity of flowshop and jobshop scheduling.
\textit{Mathematics of Operations Research}, 1976, 1(2): 117--129.

\bibitem{ortools}
Perron L, Didier F. CP-SAT. Google OR-Tools, \url{https://developers.google.com/optimization}.

\end{thebibliography}

\end{document}
